% GENERATED FILE. Edit canonical lessons, then run npm run generate:books.
% canonical-source-hash: fnv1a64:d5794e18cf2873ab
% canonical-lessons: ES-C300-sed, ES-C300-sueno, ES-C300-piel, ES-C300-sangre, ES-C300-remedio

\chapter{What the Body Says --- Thirst, Sleep, Skin, Blood, Remedy}
\label{ch:el-cuerpo-dice}
\hlchaptermodality{300}

\begin{chapteropening}
\textbf{By the end of this chapter:} \emph{I can say in Spanish that I am thirsty or sleepy, name the skin and the blood, and say there is no remedy.}

Complete ES-C300-remedio: name all five, and put three of them into a sentence with tengo or no hay.
\end{chapteropening}

\section[la sed]{la sed — something you hold, not something you are}

\label{lesson:ES-C300-sed}



\begin{quote}
Before a new run: what does \emph{una buena cosecha} mean on a bottle? (A good vintage.) And what did \emph{pascere} mean? (To put out to feed.)
\end{quote}

\subsection*{What to know first}
\begin{itemize}
\raggedright
  \item beber — what this sends you to do.
  \item tener — the word Spanish uses to report it.
\end{itemize}

\begin{sounds}
\begin{itemize}
\raggedright
  \item \emph{SED}, one beat, and the final \emph{d} is barely there — closer to a soft \emph{th} than to an English \emph{d}. $\to$ reference
\end{itemize}
\end{sounds}

\begin{cousinweb}
\textbf{la sed} — \textquotedblleft{}the thirst.\textquotedblright{}

Latin \textbf{sitis}, worn short. English inherited nothing usable from it, which is rarer than it sounds — most of the body words you have met arrived with a cousin in tow, and this one comes empty-handed.

The interesting part is not where it came from but how Spanish uses it. English says \emph{I am thirsty}, as though thirst were a description of you. Spanish says \textbf{tengo sed} — \textquotedblleft{}I have thirst,\textquotedblright{} as though it were a thing in your hands that you could put down.

Once you notice, it is everywhere: \emph{tengo hambre}, \emph{tengo frío}, \emph{tengo miedo}. Spanish keeps handing you the sensation as an object. English keeps turning it into a description of the person, which quietly makes it more permanent than it is.

Neither is right. But if you say \emph{estoy sed}, a Spanish speaker will hear something ungrammatical, and this is the reason.
\end{cousinweb}

\begin{grammarlens}[title={What you've built}]
One word, and a habit that will serve you for a dozen more.

When you want to report how your body is, reach for \emph{tengo} first and look for the thing you are holding. That is the Spanish shape.
\end{grammarlens}

\subsection*{Your turn}
Say these aloud:

\begin{itemize}
\raggedright
  \item \emph{la sed}
  \item \emph{tengo sed}
  \item \emph{tengo sed, agua por favor}
\end{itemize}

\begin{tcolorbox}[breakable,colback=teal!4,colframe=teal!35!black,title={Before you move on}]
What was \emph{sitis}? (Thirst.) Does Spanish say you are thirsty or that you have thirst? (That you have it.) And what is the whole sentence? (\emph{Tengo sed}.)
\end{tcolorbox}
\section[el sueño]{el sueño — one word for two things}

\label{lesson:ES-C300-sueno}



\begin{quote}
Two back: what did \emph{pascere} mean? (To put out to feed.) And does Spanish say you are thirsty, or that you have thirst? (That you have it.)
\end{quote}

\subsection*{What to know first}
\begin{itemize}
\raggedright
  \item dormir — the doing; this is the having.
  \item la cama — where both happen.
\end{itemize}

\begin{sounds}
\begin{itemize}
\raggedright
  \item \emph{SUE-ño}, stress on the first part. The \emph{ue} is one slide, and the \emph{ñ} is the \emph{n} with a \emph{y} leaning on it. $\to$ reference
\end{itemize}
\end{sounds}

\begin{cousinweb}
\textbf{el sueño} — \textquotedblleft{}sleep,\textquotedblright{} and also \textquotedblleft{}a dream.\textquotedblright{}

Latin \textbf{somnus}, sleep. The stressed \emph{o} cracked open into \emph{ue}, exactly as it did turning \emph{dormir} into \emph{duermo} and \emph{porta} into \emph{puerta}. Once that crack is familiar you can run it backwards: see \emph{ue} in a Spanish word, put an \emph{o} back, and look for the Latin.

Do it here and \emph{somnus} appears — and with it \textbf{insomnia}, which is \emph{in-}, not, plus sleep; \textbf{somnolent}, heavy with it; and a \textbf{somnambulist}, who walks in it.

The double meaning is worth a moment. \emph{El sueño} is the sleep and the dream at once, and Spanish is content to leave them together. \emph{Tengo sueño} is \textquotedblleft{}I am sleepy\textquotedblright{} — you are holding sleep, in the way you hold thirst. \emph{Tuve un sueño} is \textquotedblleft{}I had a dream.\textquotedblright{} Same word, and the difference sits in everything around it.
\end{cousinweb}

\begin{grammarlens}[title={What you've built}]
Two of the body, both of them things you hold rather than things you are.

You have had \emph{dormir} for the act since early on. \emph{El sueño} is what the act is made of, and what it leaves behind in the morning.
\end{grammarlens}

\subsection*{Your turn}
Say these aloud:

\begin{itemize}
\raggedright
  \item \emph{el sueño}
  \item \emph{tengo sueño}
  \item \emph{tengo sed}, then \emph{tengo sueño}
\end{itemize}

\begin{tcolorbox}[breakable,colback=teal!4,colframe=teal!35!black,title={Before you move on}]
What was \emph{somnus}? (Sleep.) What happened to its stressed \emph{o}? (It broke into \emph{ue}.) And what two things does \emph{el sueño} name? (Sleep and a dream.)
\end{tcolorbox}
\section[la piel]{la piel — the hide you are wearing}

\label{lesson:ES-C300-piel}



\begin{quote}
What is the Spanish shape for thirst? (\emph{Tengo sed} — you have it.) And what happened to the \emph{o} of \emph{somnus}? (It broke into \emph{ue}.)
\end{quote}

\subsection*{What to know first}
\begin{itemize}
\raggedright
  \item la mano — one of the places you can see it best.
  \item la oveja — an animal people keep partly for hers.
\end{itemize}

\begin{sounds}
\begin{itemize}
\raggedright
  \item \emph{PIEL}, one beat. The \emph{ie} is a single slide and the \emph{l} stays bright at the end. $\to$ reference
\end{itemize}
\end{sounds}

\begin{cousinweb}
\textbf{la piel} — \textquotedblleft{}the skin.\textquotedblright{}

Latin \textbf{pellis} was a hide — an animal's, taken off. The stressed \emph{e} broke into \emph{ie}, the same crack you have watched in \emph{cielo} and \emph{miel}, and the word came out meaning the skin you are still standing in.

That is worth sitting with. Spanish uses one word for the skin on your arm and the skin a coat is made of. \emph{Una chaqueta de piel} is a leather jacket. Latin drew no line there either, and English mostly does: \textbf{skin} for the living one, \textbf{hide} and \textbf{pelt} for the dead.

\textbf{Pelt} is \emph{pellis} itself, arrived through French. And the strangest cousin is in a church: a \textbf{surplice} is \emph{super-pellicium}, the white linen worn \textbf{over} the fur coat, because medieval choirs were cold. The garment is named entirely for what is underneath it.
\end{cousinweb}

\begin{grammarlens}[title={What you've built}]
Three of the body, and the second one whose \emph{ie} you could have predicted.

\emph{Cielo}, \emph{miel}, \emph{piel}: three words, one crack, and it always runs the same way. Put an \emph{e} back where the \emph{ie} is and the Latin will usually show itself.
\end{grammarlens}

\subsection*{Your turn}
Say these aloud:

\begin{itemize}
\raggedright
  \item \emph{la piel}
  \item \emph{la piel de la mano}
  \item \emph{el sueño}, then \emph{la piel}
\end{itemize}

\begin{tcolorbox}[breakable,colback=teal!4,colframe=teal!35!black,title={Before you move on}]
What was \emph{pellis}? (A hide.) Which English word is it, almost untouched? (\emph{Pelt}.) And what is a surplice worn over? (A fur coat.)
\end{tcolorbox}
\section[la sangre]{la sangre — a word that used to be a diagnosis}

\label{lesson:ES-C300-sangre}



\begin{quote}
Three questions. What was \emph{sitis}? (Thirst.) What was \emph{somnus}? (Sleep.) And what was \emph{pellis}? (A hide.)
\end{quote}

\subsection*{What to know first}
\begin{itemize}
\raggedright
  \item el vaso — the word for a vessel, which the body has its own set of.
  \item la piel — what keeps this one where it belongs.
\end{itemize}

\begin{sounds}
\begin{itemize}
\raggedright
  \item \emph{SAN-gre}, stress on the first part. The \emph{g} is hard and the \emph{r} is a single tap. $\to$ reference
\end{itemize}
\end{sounds}

\begin{cousinweb}
\textbf{la sangre} — \textquotedblleft{}the blood.\textquotedblright{}

Latin \textbf{sanguis}, and through its sideways form \textbf{sanguinem}, which is where the \emph{-re} on the end comes from. The road is short and the meaning never moved.

English took the same word and did something odd with it. Medieval medicine held that a person's temper came from whichever of four fluids ran strongest in them, and somebody with plenty of blood was ruddy, warm and hopeful. That is why \textbf{sanguine} still means optimistic — a medical theory abandoned four hundred years ago, fossilised in an ordinary English word.

The plainer cousins are easier. \textbf{Consanguinity} is shared blood, which is what a family is. And French gave English \textbf{sangfroid}, \textquotedblleft{}cold blood\textquotedblright{}: the calm of somebody whose blood has not risen.

When you met \emph{el vaso}, you met the word for a vessel, and the body's own vessels are what this runs through. The two words have been waiting for each other.
\end{cousinweb}

\begin{grammarlens}[title={What you've built}]
Four of the body: two you hold, one you wear, one that runs underneath.

\emph{La sed} and \emph{el sueño} are sensations. \emph{La piel} and \emph{la sangre} are the body itself, and between them they are the outside and the inside.
\end{grammarlens}

\subsection*{Your turn}
Say these aloud:

\begin{itemize}
\raggedright
  \item \emph{la sangre}
  \item \emph{la piel}, then \emph{la sangre}
  \item all four — \emph{la sed}, \emph{el sueño}, \emph{la piel}, \emph{la sangre}
\end{itemize}

\begin{tcolorbox}[breakable,colback=teal!4,colframe=teal!35!black,title={Before you move on}]
What was \emph{sanguis}? (Blood.) Why does \emph{sanguine} mean cheerful? (An old theory that blood made you ruddy and hopeful.) And what is \emph{sangfroid}, taken literally? (Cold blood.)
\end{tcolorbox}
\section[el remedio]{el remedio — healing, taken again}

\label{lesson:ES-C300-remedio}



\begin{quote}
Before the last of the five: what was \emph{sitis}? (Thirst.) What broke into \emph{ue} in \emph{sueño}? (The stressed \emph{o} of \emph{somnus}.) And why does \emph{sanguine} mean cheerful? (An old theory about blood.)
\end{quote}

\subsection*{What to know first}
\begin{itemize}
\raggedright
  \item la fiebre — one of the things this is for.
  \item la enfermera — who is likely to bring it.
\end{itemize}

\begin{sounds}
\begin{itemize}
\raggedright
  \item \emph{re-ME-dio}, stress in the middle. The \emph{r} at the start is a single tap, and the \emph{io} at the end is one slide. $\to$ reference
\end{itemize}
\end{sounds}

\begin{cousinweb}
\textbf{el remedio} — \textquotedblleft{}the remedy,\textquotedblright{} the cure, the way out of a trouble.

Two pieces. \textbf{Re-}, again or back, which you have seen doing this work before. And \textbf{mederi}, Latin for to heal or to look after. A \emph{remedium} is a putting-back of health.

\emph{Mederi} is one of the more productive roots English has. \textbf{Medicine} is the healing art. A \textbf{medic} is the person practising it. Something \textbf{remedial} is meant to put a thing back. All of them come off the same short word.

The phrase to take away is \textbf{no hay remedio} — \textquotedblleft{}there is no remedy,\textquotedblright{} said about a thing that cannot be helped. You already have \emph{hay} and you already have \emph{no}, so the sentence is finished the moment you have this word. It is what a Spanish speaker says with a shrug where English says \emph{it can't be helped}, and it is useful far outside a sickroom.
\end{cousinweb}

\begin{grammarlens}[title={What you've built}]
Five: \emph{la sed}, \emph{el sueño}, \emph{la piel}, \emph{la sangre}, \emph{el remedio}.

Read in order they are a body reporting on itself and then being attended to: two things it wants, two things it is made of, one thing that mends it.

And four of the five you can put into a sentence with words you already had — \emph{tengo sed}, \emph{tengo sueño}, \emph{no hay remedio}. That is the difference between a list of words and a run you can speak from.
\end{grammarlens}

\subsection*{Your turn}
Say these aloud:

\begin{itemize}
\raggedright
  \item \emph{el remedio}
  \item \emph{no hay remedio}
  \item all five — \emph{la sed}, \emph{el sueño}, \emph{la piel}, \emph{la sangre}, \emph{el remedio}
\end{itemize}

\begin{tcolorbox}[breakable,colback=teal!4,colframe=teal!35!black,title={Before you move on}]
What does \emph{mederi} mean? (To heal.) Which everyday English word is the same root? (\emph{Medicine}.) And what does \emph{no hay remedio} say? (There is nothing to be done.)
\end{tcolorbox}
